\section{Reduction 2-Complexes and Elementary Holonomy}
\label{sec:reduction-two-complex}

This section introduces the object on which the flatness reading of confluence is literal. It adds to the reduction graph a list of \emph{branchings}, the pairs of one-step rewrites to be compared, and fills a branching with a $2$-cell exactly when its two branches can be brought together.

\subsection{Branchings and fillers}

\begin{definition}[Elementary peaks and fillers]
\label{def:elementary-peak-flatness}
\label{def:boundary-filler-data}
Let $(X,\to)$ be an abstract rewriting system.
\begin{enumerate}
    \item An \emph{elementary peak} is a triple $(a,b,c)$ with $a \to b$, $a \to c$, and $b \neq c$. We call $a$ its \emph{source} and $b$, $c$ its \emph{branches}. The triples $(a,b,c)$ and $(a,c,b)$ describe the same \emph{branching}; when branchings are counted or cells attached, the two are identified and one representative is used.
    \item A \emph{filler} for the peak $(a,b,c)$ is a state $w$, the \emph{join witness}, together with rewrite sequences $b \to^{*} w$ and $c \to^{*} w$.
\end{enumerate}
\end{definition}

A filler turns the open shape $b \leftarrow a \to c$ into a closed walk
\[
a \to b \to^{*} w \leftarrow^{*} c \leftarrow a,
\]
which can serve as the boundary of a disk. Without a filler there is no closed boundary to attach anything to. This asymmetry is the whole point of the construction, and we build it into the definition.

\begin{definition}[Reduction 2-complex]
\label{def:reduction-two-complex}
Let $(X,\to)$ be an abstract rewriting system and let $G$ be a set of branchings, i.e.\ of elementary peaks up to exchange of the two branches, called the \emph{generating branchings}. The \emph{reduction 2-complex} of $(X,\to)$ with respect to $G$ consists of
\begin{enumerate}
    \item the states of $X$ as $0$-cells;
    \item the one-step rewrites $a \to b$ as directed $1$-cells;
    \item for each branching in $G$ that admits a filler, one $2$-cell attached along the closed walk of a chosen filler.
\end{enumerate}
Branchings in $G$ that admit no filler are kept as \emph{open branchings}: they are recorded, but no cell is attached to them.
\end{definition}

For an abstract rewriting system, $G$ is the set of all branchings. For string and term rewriting, \Cref{sec:critical-pairs-curvature} takes $G$ to be the set of critical-pair instances, the branchings that can actually fail. Whether a branching admits a filler does not depend on which filler is chosen; the choice affects only which $2$-cell is drawn. The computations of \Cref{sec:computational-bridge} fix a canonical filler by a deterministic rule (\Cref{app:computational-protocol}) so that the recorded complexes are reproducible.

\subsection{Elementary holonomy}

\begin{definition}[Elementary holonomy and elementary flatness]
\label{def:elementary-holonomy-status}
An elementary peak has \emph{trivial elementary holonomy} if it admits a filler, and \emph{nontrivial elementary holonomy} otherwise; the status is the same for $(a,b,c)$ and $(a,c,b)$, so it is a property of the branching. A rewriting system is \emph{elementarily flat} if every elementary peak has trivial elementary holonomy, that is, if the branches of every elementary peak are joinable. A reduction 2-complex records this status for the branchings in $G$: the filled ones carry $2$-cells and the others remain open.
\end{definition}

The term is chosen to evoke, not to import, the geometric notion. Elementary holonomy is a yes/no status attached to a single branching: either going down one branch and then on to $w$ agrees with going down the other branch and then on to $w$, for some $w$, or no such $w$ exists. It is not the holonomy of a connection valued in a group, and nothing below is a statement about fundamental groups, homotopy of rewrite paths, or higher coherence. Those richer questions are the subject of the homotopical theory of rewriting initiated by Squier. For a convergent (terminating and confluent) presentation of a monoid, Squier showed that a choice of confluence diagrams for the critical branchings yields a homotopy basis of the associated space of derivations, and the polygraphic setting extends this to higher dimensions \cite{Squier1987,SquierOtto1987,GuiraudMalbos2012,GuiraudMalbos2018,Mimram2023}. Elementary holonomy records only the first, Boolean layer of that structure: whether a given $2$-cell exists at all.

\begin{remark}[Comparing normal forms is a different test]
\label{rem:holonomy-vs-formal-core}
For terminating systems it is tempting to declare a peak defective when its two branches have different sets of reachable normal forms. This is not the same as nonjoinability, and it is not local. In the system of \Cref{fig:nf-mismatch}, the peak at $A$ is joinable at $D$, so its holonomy is trivial, yet $B$ reaches the normal forms $\{D,E\}$ while $C$ reaches only $\{D\}$. The discrepancy is caused by a different peak, the one at $B$, which is genuinely nonjoinable. In our computations, normal-form comparison is therefore recorded only as a secondary diagnostic, and only for terminating systems; elementary holonomy and all defect counts use joinability alone.
\end{remark}

\begin{figure}[t]
\centering
\begin{tikzpicture}
  \node[st] (A) at (0,0) {$A$};
  \node[st] (B) at (-1.1,-1.1) {$B$};
  \node[st] (C) at (1.1,-1.1) {$C$};
  \node[nf] (D) at (0,-2.2) {$D$};
  \node[nf] (E) at (-2.2,-2.2) {$E$};
  \begin{scope}[on background layer]
    \fill[cellfill] (A.center) -- (B.center) -- (D.center) -- (C.center) -- cycle;
  \end{scope}
  \draw[rw] (A)--(B); \draw[rw] (A)--(C);
  \draw[rw] (B)--(D); \draw[rw] (C)--(D); \draw[rw] (B)--(E);
  \draw[defect] (E) .. controls (-1.4,-2.8) and (-0.6,-2.8) .. (D);
  \node[font=\scriptsize, align=left, anchor=west] at (2.0,-0.6) {peak at $A$: joinable at $D$ (trivial)};
  \node[font=\scriptsize, align=left, anchor=west] at (2.0,-1.1) {normal forms: $B \mapsto \{D,E\}$, $C \mapsto \{D\}$};
  \node[font=\scriptsize, align=left, anchor=west] at (2.0,-1.6) {peak at $B$: $D$, $E$ not joinable (nontrivial)};
\end{tikzpicture}
\caption{Joinability and normal-form agreement differ. The peak at $A$ is filled (shaded) because $B$ and $C$ both reach $D$, even though the sets of normal forms reachable from $B$ and from $C$ differ. The obstruction to confluence sits at the peak $D \leftarrow B \to E$.}
\label{fig:nf-mismatch}
\end{figure}

\subsection{Representative examples}

The same local question arises in each of the settings studied later: given two one-step rewrites from a common state, is there a filler? \Cref{tab:elementary-holonomy-summary} lists the generating branchings and their holonomy in six representative examples. In the finite systems $E001$ and $E002$, the generating branchings are the peaks of \Cref{fig:bare-graph-failure}. In the string pair $E007$--$E008$ and the term pair $TRS003$--$TRS004$ they are critical-pair instances; these pairs are discussed in \Cref{sec:critical-pairs-curvature,sec:completion-holonomy}.

\begin{table}[t]
    \caption{Generating branchings in six representative examples. A filled branching carries a $2$-cell attached along the closed walk through its join witness; an open branching has nontrivial elementary holonomy and no attached cell.}
    \label{tab:elementary-holonomy-summary}
    \centering
    \smallskip
    \small
\begin{tabular}{@{}llcccp{0.20\linewidth}@{}}
\toprule
Example & Domain & 2-cells & Trivial & Nontrivial & Witness / defect \\
\midrule
$E001$ (diamond)       & finite ARS & 1 & 1 & 0 & D \\
$E002$ (fork)          & finite ARS & 1 & 0 & 1 & nonjoinable \\
$E007$ (before)        & string     & 1 & 0 & 1 & nonjoinable \\
$E008$ (after)         & string     & 1 & 1 & 0 & \texttt{y} \\
$TRS003$ (nested)      & term rew.  & 1 & 0 & 1 & nonjoinable \\
$TRS004$ (resolved)    & term rew.  & 1 & 1 & 0 & \texttt{p(a)} \\
\bottomrule
\end{tabular}

\end{table}
